\subsection{محدودیت‌ها و تحلیل شکست}
بررسی فرآیند فشرده‌سازی و خطاهای رخ‌داده بیانگر نکات تحلیلی زیر است:
\begin{itemize}
    \item \textbf{افت عملکرد نسبی در مدل \lr{Qwen} (دقت \lr{51.99\%}):} در مدل زبانی \lr{Qwen} دقت برابر با \lr{51.99\%} حاصل شده است. علت اصلی این پدیده، وجود وزن‌های متمرکز با مقادیر پرت بسیار شدید در لایه‌های چندگروهی (\lr{GQA}) مدل‌های مدرن است؛ پیش‌شرط‌سازی ساده با گشتاور دوم در غیاب داده‌های کالیبراسیون غنی و متنوع زبانی، نتوانسته جهت‌گیری تمام کانال‌های پرت را به طور کامل حفظ کند و منجر به تخریب نسبی فاکتورهای کم‌مرتبه در لایه‌های انتهایی دیکودر شده است.
    \item \textbf{حساسیت به فرآیند کالیبراسیون و انحراف لایه‌نرمال:} در صورتی که متن‌های کالیبراسیون ورودی توزیع ناکافی داشته باشند یا پارامتر \lr{recalibrate\_ln} غیرفعال گردد، انحراف واریانس فعال‌سازی در طول لایه‌های عمیق تشدید شده و موجب افزایش خطای فروبنیوس فراتر از حد مجاز (\lr{max\_error}) می‌شود که در آن صورت ماژول ناگزیر به استفاده از مسیر پشتیبان (\lr{fallback}) و بازگردانی لایه به حالت متراکم می‌گردد.
\end{itemize}
